\documentclass[11pt]{article}
\usepackage[margin=1in]{geometry}
\usepackage[T1]{fontenc}
\usepackage[utf8]{inputenc}
\usepackage{lmodern}
\usepackage{amsmath,amssymb,amsthm,mathtools}
\usepackage{enumitem}
\usepackage{booktabs,longtable,array}
\usepackage[hidelinks]{hyperref}

\newtheorem{definition}{Definition}[section]
\newtheorem{theorem}[definition]{Theorem}
\newtheorem{proposition}[definition]{Proposition}
\newtheorem{lemma}[definition]{Lemma}
\newtheorem{corollary}[definition]{Corollary}
\newtheorem{assumption}[definition]{Assumption}
\newtheorem{remark}[definition]{Remark}

\newcommand{\Id}{\operatorname{Id}}
\newcommand{\Hist}{\operatorname{Hist}}
\newcommand{\Resp}{\operatorname{Resp}}
\newcommand{\Adm}{\operatorname{Adm}}
\newcommand{\QGauge}{\operatorname{Gauge}}
\newcommand{\Aut}{\operatorname{Aut}}
\newcommand{\Iint}{I_{\mathrm{int}}}
\newcommand{\ThetaS}{\Theta_S}
\newcommand{\Aset}{\mathcal{A}}
\newcommand{\Fset}{\mathcal{F}}
\newcommand{\R}{\mathbb{R}}

\title{QICN v18 Internal Note: Product-Separator Counterexample and Conditional Closure of $I_{\mathrm{int}}$}
\author{QICN internal formal-audit scaffold}
\date{2026-05-26}

\begin{document}
\maketitle

\section*{Governance boundary}
This document does not certify external support, consciousness, phenomenality, identity transfer, or bridge-burden closure. It is an internal mathematical scaffold. Its assertions are limited to the formal definitions and conditional results stated below. It does not report human-subject data, external adjudication, reviewer certification, or empirical validation.

\section{Scope and separation of layers}
The purpose of this note is to close the internal $\Iint$ gap as far as the present formal resources allow, without inflating the public status of the framework.

\begin{description}[leftmargin=2.6cm]
  \item[Ontology.] A system $S$ is treated as an admissible object with identity data, histories, interventions, and response observations. No consciousness, agency, subjecthood, or phenomenal identity is attributed by this note.
  \item[Mathematical model.] The model is a category-like typed setting with admissible systems, admissible factorizations, separator families, and gauge transformations preserving the declared data.
  \item[Implementation.] No source paper is edited by this note. The note can be compiled independently and linked from reports or registries.
  \item[Language.] Terms such as ``atomic'', ``separator'', and ``integration'' are formal predicates or functionals defined below. Any intuitive reading is non-operative unless it reduces to these predicates.
  \item[Interpretation.] The strongest closed result below is conditional: $\Iint$ is closed under an explicit atomic or prime-separator hypothesis. The product-separator construction proves that the original weaker assumptions do not imply atomicity by themselves.
\end{description}

\section{Typed setting}

\begin{definition}[Admissible system package]
An admissible system package is a tuple
\[
  S = (X,\mathcal{U},\Hist_S,\Id_S,\Resp_S,\Theta_S),
\]
where $X$ is a state space, $\mathcal{U}$ is a set of admissible interventions, $\Hist_S$ is the admissible history family, $\Id_S$ is the rigid identity datum, $\Resp_S$ is the intervention-response assignment, and $\Theta_S$ is a family of separators used to distinguish identity/history/response alternatives.
\end{definition}

\begin{definition}[Exact admissible factorization]
An exact admissible factorization of $S$ is a tuple
\[
  F:S\longrightarrow S_1\times S_2
\]
with non-terminal factors $S_1,S_2$ such that the state, intervention, history, identity, and response data decompose as
\[
  X\cong X_1\times X_2,
  \quad
  \mathcal U\cong \mathcal U_1\times\mathcal U_2,
  \quad
  \Hist_S\cong \Hist_{S_1}\times\Hist_{S_2},
\]
\[
  \Id_S \cong (\Id_{S_1},\Id_{S_2}),
  \quad
  \Resp_S(u_1,u_2,x_1,x_2)=(\Resp_{S_1}(u_1,x_1),\Resp_{S_2}(u_2,x_2)).
\]
The factorization is non-trivial when neither factor is isomorphic to the terminal admissible system.
\end{definition}

\begin{definition}[Factor-local separator partition]
A separator partition for a factorization $F:S\to S_1\times S_2$ is a disjoint union
\[
  \Theta_S = \Theta_1\sqcup\Theta_2
\]
such that every $\theta\in\Theta_i$ is pulled back from factor $S_i$ and separates only identity/history/response alternatives visible in that factor.
\end{definition}

\begin{definition}[Atomic separator]
The separator family $\Theta_S$ is atomic for $S$ when there is no exact non-trivial admissible factorization $F:S\to S_1\times S_2$ admitting a factor-local separator partition that preserves $\Id_S$, $\Hist_S$, and $\Resp_S$ simultaneously.
\end{definition}

\begin{remark}[Non-circularity convention]
The phrase ``atomic separator'' is a formal non-decomposition predicate. A later axiom that merely restates this definition under another name would be circular. A useful additional condition must instead be stated in terms of checkable intervention support, response coupling, commutants, graph connectivity, or perturbation behavior, and only then shown to imply atomicity.
\end{remark}

\section{Product-separator counterexample}

\begin{proposition}[Product separator blocks the global derivation]
Rigid identity, product-topological continuity, and componentwise intervention fidelity do not imply atomicity of $\Theta_S$.
\end{proposition}

\begin{proof}
Let $S_1$ and $S_2$ be two admissible system packages and form the product package
\[
  S = S_1\times S_2.
\]
Set
\[
  X=X_1\times X_2,
  \quad
  \mathcal U=\mathcal U_1\times\mathcal U_2,
  \quad
  \Hist_S=\Hist_{S_1}\times\Hist_{S_2},
\]
\[
  \Id_S=(\Id_{S_1},\Id_{S_2}),
  \quad
  \Resp_S((u_1,u_2),(x_1,x_2))=(\Resp_{S_1}(u_1,x_1),\Resp_{S_2}(u_2,x_2)).
\]
Give $X$ the product topology. If the component response maps are continuous, the product response map is continuous. If each component is intervention-faithful relative to its own data, the product is faithful componentwise: equality of product responses implies equality in each coordinate at the relevant observational resolution.

Now take
\[
  \Theta_S = \pi_1^*\Theta_{S_1}\sqcup \pi_2^*\Theta_{S_2},
\]
where $\pi_i$ are the coordinate projections. This separator family preserves identity, histories, and responses by construction, but it decomposes into factor-local separator families. Therefore the product factorization is non-trivial and admits a factor-local separator partition. Hence $\Theta_S$ is not atomic.

The construction satisfies the upstream rigid-identity, continuity, and componentwise fidelity requirements, yet violates separator atomicity. Thus the global implication from the weaker assumptions to atomicity is false in this model class.
\end{proof}

\begin{corollary}[Exact remaining burden]
Any proof of a global $I_{\mathrm{int}}$ theorem must add a condition that excludes product-separator systems or must restrict the admissible system class so that the product construction is inadmissible.
\end{corollary}

\section{Candidate strengthening conditions}

\begin{definition}[No-cloning of intervention channels]
An admissible system satisfies intervention no-cloning when no non-trivial factorization $S\cong S_1\times S_2$ admits two factor-local intervention subfamilies $\mathcal U_1,\mathcal U_2$ whose pulled-back response maps reproduce the full response assignment while preserving $\Id_S$ and $\Hist_S$.
\end{definition}

\begin{definition}[Connected response-support graph]
Let $G_S$ be the bipartite graph whose left vertices are elementary separator tests, whose right vertices are elementary intervention-response coordinates, and whose edges indicate that a separator test changes under a perturbation of that coordinate within the admissible response model. The system has connected response support when $G_S$ is connected after deleting isolated non-operative vertices.
\end{definition}

\begin{definition}[Kernel robustness under local perturbations]
Let $D\Resp_S$ be the linearized response operator in a chosen admissible chart. The system has locally robust kernel coupling when every sufficiently small factor-local perturbation preserving $\ker(D\Resp_S)$ also preserves the full coupled response-support graph. Equivalently, there is no local perturbation that splits $\ker(D\Resp_S)$ into two invariant factor kernels while preserving $\Id_S$ and $\Hist_S$.
\end{definition}

\begin{definition}[Prime intervention-response coupling]
A system has prime intervention-response coupling when, for every exact non-trivial admissible factorization candidate $F:S\to S_1\times S_2$, at least one of the following fails:
\begin{enumerate}[label=(\roman*)]
  \item response support decomposes as a disjoint union of factor supports;
  \item the separator family is the pullback union of factor-local separators;
  \item the gauge group preserving identity, histories, and responses splits as a product of factor gauge groups;
  \item the linearized response operator has an invariant block decomposition compatible with the same factorization.
\end{enumerate}
\end{definition}

\begin{lemma}[Connected support blocks product separators]
If $S$ has connected response support and every separator in $\Theta_S$ is response-operative, then no product-separator partition can preserve the full response support.
\end{lemma}

\begin{proof}
Assume a product-separator partition $\Theta_S=\Theta_1\sqcup\Theta_2$ exists for a non-trivial factorization. Since each separator is factor-local, perturbations of response coordinates in factor $S_1$ can affect only tests in $\Theta_1$, and perturbations in factor $S_2$ can affect only tests in $\Theta_2$. The response-support graph then decomposes into the disjoint union of the $S_1$ and $S_2$ support graphs. Because all separators are response-operative, both subgraphs contain at least one edge. This contradicts connectedness of $G_S$.
\end{proof}

\begin{lemma}[Prime coupling implies separator atomicity]
If $S$ has prime intervention-response coupling, then $\Theta_S$ is atomic.
\end{lemma}

\begin{proof}
Suppose $\Theta_S$ is not atomic. Then by definition there exists an exact non-trivial admissible factorization admitting a factor-local separator partition preserving identity, histories, and responses. Such a factorization satisfies separator pullback decomposition. Because the factorization is exact and intervention-faithful, it also gives a response-support decomposition, a compatible split gauge action, and an invariant block decomposition of the linearized response operator. This contradicts prime intervention-response coupling, which requires at least one of these four clauses to fail for every non-trivial factorization candidate.
\end{proof}

\begin{longtable}{p{0.22\textwidth}p{0.22\textwidth}p{0.22\textwidth}p{0.22\textwidth}}
\toprule
Candidate & Blocks product separator? & Risk of circularity & Audit status \\
\midrule
No-cloning of intervention channels & Yes, if stated over response reproduction rather than separator names. & Medium. Too strong if it simply says no duplicated separator. & Useful but needs operational tests. \\
Connected response-support graph & Yes, under response-operative separator assumption. & Low to medium. It is graph-theoretic and testable on a model. & Strongest non-verbal candidate. \\
Kernel robustness & Yes, if factor kernels are excluded under perturbation. & Medium. Depends on chart and linearization. & Good analytic supplement, not sufficient alone. \\
Prime intervention-response coupling & Yes. & Medium-high. It is close to atomicity, but stated over response/gauge/operator decomposition rather than direct separator naming. & Best aggressive closure axiom for v18, still conditional. \\
\bottomrule
\end{longtable}

\section{Conditional closure of $I_{\mathrm{int}}$}

\begin{definition}[$I_{\mathrm{int}}$ as a factorization-gap functional]
For an admissible system $S$, define the exact factorization gap
\[
  \Delta_{\mathrm{int}}(S)=\inf_{F\in\mathcal F_{\mathrm{nt}}(S)}
  \Bigl[\mathcal L_{\mathrm{hist}}(F)+\mathcal L_{\mathrm{id}}(F)+\mathcal L_{\mathrm{resp}}(F)+\mathcal P_{\mathrm{gauge}}(F)\Bigr],
\]
where $\mathcal F_{\mathrm{nt}}(S)$ is the class of non-trivial admissible factorizations, the loss terms vanish exactly when histories, identity, and responses are preserved, and $\mathcal P_{\mathrm{gauge}}$ vanishes exactly when the factorization is compatible with the admissible gauge action. Then
\[
  \Iint(S)=1 \quad\Longleftrightarrow\quad \Delta_{\mathrm{int}}(S)>0.
\]
If $\mathcal F_{\mathrm{nt}}(S)$ is empty, set $\Delta_{\mathrm{int}}(S)=+\infty$.
\end{definition}

\begin{theorem}[Conditional on $\Theta_S$ atomic: existence of $I_{\mathrm{int}}$]
If $\Theta_S$ is atomic and the loss terms above are faithful to history, identity, response, and gauge preservation, then $\Iint(S)$ exists as a well-defined binary invariant of the admissible isomorphism class of $S$.
\end{theorem}

\begin{proof}
The functional $\Delta_{\mathrm{int}}$ is defined as an infimum over a class determined by the admissible isomorphism class of $S$. The loss terms are also invariant under admissible isomorphism because they are expressed only in terms of preserved histories, identity, responses, and gauge compatibility. Thus $\Delta_{\mathrm{int}}(S)$ is well-defined on admissible isomorphism classes.

Atomicity of $\Theta_S$ rules out exact non-trivial factorizations with zero loss. Faithfulness of the loss terms implies that zero loss is equivalent to preserving the four relevant structures. Therefore the zero-loss case is impossible for non-trivial factorizations. The predicate $\Delta_{\mathrm{int}}(S)>0$ is consequently a well-defined formal invariant whenever a positive margin is part of the admissibility criterion. The binary value $\Iint(S)$ is that predicate.
\end{proof}

\begin{theorem}[Conditional on $\Theta_S$ atomic: uniqueness under gauge]
Assume $\Theta_S$ is atomic and the admissible gauge group $\QGauge(S)$ acts by preserving histories, identity, responses, and separator membership. Then any binary invariant $J(S)$ that is zero exactly on systems admitting a zero-loss non-trivial factorization and is invariant under $\QGauge(S)$ agrees with $\Iint(S)$.
\end{theorem}

\begin{proof}
By definition, $\Iint(S)=0$ exactly when $\Delta_{\mathrm{int}}(S)=0$, which is exactly the existence of a zero-loss non-trivial factorization. The invariant $J$ is stipulated to have the same zero-set. Both $J$ and $\Iint$ are binary. Therefore they agree pointwise. Gauge invariance ensures that the value is not changed by reparameterizations of the separator or response presentation.
\end{proof}

\begin{corollary}[Aggressive but bounded v18 closure]
Under prime intervention-response coupling, response-operative separators, and faithful factorization losses, $\Iint$ is conditionally closed as a unique gauge-invariant binary factorization-gap invariant.
\end{corollary}

\begin{proof}
Prime intervention-response coupling implies separator atomicity. The two conditional theorems then give existence and uniqueness of $\Iint$ as a gauge-invariant binary invariant.
\end{proof}

\section{What is closed and what remains open}

\begin{longtable}{p{0.25\textwidth}p{0.22\textwidth}p{0.43\textwidth}}
\toprule
Item & Status & Reason \\
\midrule
Product-separator counterexample & Closed in this note & It explicitly satisfies the weaker assumptions and fails atomicity. \\
Global derivation of atomicity from rigid identity + continuity + fidelity & Not closed & The counterexample blocks the implication. \\
$\Iint$ under atomic separator & Conditionally closed & Existence and uniqueness follow from atomicity plus faithful factorization losses. \\
$\Iint$ under prime intervention-response coupling & Conditionally closed & Prime coupling implies atomicity, then the conditional theorem applies. \\
Empirical or external support & Not claimed & No external campaign, human review, or adjudication is executed here. \\
\bottomrule
\end{longtable}

\section{Implementation note}
This note is intentionally not inserted into the original Paper 5 source. It is a v18 scaffold that can be reviewed, criticized, or promoted by a later human mathematical curation process. Promotion into the theorem registry would require explicit registry overlay, human mathematical review, and a decision about whether prime intervention-response coupling is an acceptable new hypothesis rather than a disguised restatement of separator atomicity.

\end{document}
